% ===== Wrap-up and the end of the series: 5 ideas, try at home, the afterword, the whole book, goodbye, questions (~7 min) =====
\partaccent{butter}\notopic

\begin{frame}{Այսօրվա 5 գաղափարը}
  \begin{enumerate}
    \item<1-> «Նորմալի» ու «խանգարման» սահմանը \textbf{մշուշոտ} է, իսկ հոգեկան խնդիրները շատ տարածված են
    \item<2-> Դեպրեսիան տխրություն չէ. այն փոխում է \textbf{մտածողությունը}, ու «եսասեր», «ծույլ» պիտակները սխալ են
    \item<3-> Սթրեսը վերահսկողության կորուստ է. քրոնիկ սթրեսը կարող է հասցնել \textbf{վթարային անջատման}
    \item<4-> Թմրանյութերը կարճ ճանապարհով են միացնում պարգևի ուղին, իսկ ուղեղը հարմարվում է. \textbf{տոլերանտություն} ու կախվածություն
    \item<5-> Հալյուցինացիան ներքինն ու արտաքինը \textbf{շփոթելն} է, զառանցանքը՝ անհամապատասխանության \textbf{սխալ բացատրությունը}
  \end{enumerate}
\end{frame}

\begin{frame}{Փորձիր տանը}
  \begin{columns}[T]
    \begin{column}{0.31\textwidth}
      \begin{card}[colback=mint!30!cardwhite]
        \centering\bigemoji[1cm]{feather}\par\smallskip
        \raggedright Թող ընկերդ խուտուտ տա քեզ, հետո ինքդ փորձիր նույն տեղում։ Տարբերությունը զգու՞մ ես
      \end{card}
    \end{column}
    \begin{column}{0.31\textwidth}
      \begin{card}[colback=lavender!35!cardwhite]
        \centering\bigemoji[1cm]{thought}\par\smallskip
        \raggedright Մեկ օր նկատիր, քանի տարօրինակ բացատրություն է ուղեղդ առաջարկում ու անմիջապես մերժում (ինչպես ավտոբուսի դեպքում)
      \end{card}
    \end{column}
    \begin{column}{0.31\textwidth}
      \begin{card}[colback=peach!35!cardwhite]
        \centering\bigemoji[1cm]{batterylow}\par\smallskip
        \raggedright Հաջորդ սթրեսի ժամանակ նկատիր, թե ինչին ես ձեռք մեկնում «հանգստանալու» համար։ Սա ցիկլի սկի՞զբ է
      \end{card}
    \end{column}
  \end{columns}
\end{frame}

% The afterword: the whole thing is one sentence
\begin{frame}[plain]
  \begin{tikzpicture}[remember picture, overlay]
    \fill[brainpink!45!cream] (current page.south west) rectangle (current page.north east);
    \fill[pattern={Dots[distance=4.4mm, radius=0.28mm]}, pattern color=plum!12!brainpink!45!cream]
      (current page.south west) rectangle (current page.north east);
    \fill[cardwhite, opacity=0.45] ([xshift=-3.3cm, yshift=-0.2cm]current page.east) circle (3.1cm);
    \node[anchor=north west, inner sep=0] at ([xshift=10mm, yshift=-12mm]current page.north west)
      {\idiotpill{book}{Վերջաբան}};
    \node[anchor=west, text=plum, align=left, text width=9.6cm, inner sep=0, font=\blackfont\fontsize{22}{28}\selectfont]
      (t) at ([xshift=10mm, yshift=6mm]current page.west)
      {Ահա և ուղեղը։\\ Տպավորիչ է, չէ՞։\\ Բայց նաև՝ մի քիչ ապուշ։};
    \node[anchor=north west, text=plum!75!cream, text width=9.4cm, inner sep=0, align=left, font=\mediumfont\normalsize]
      at ([yshift=-5mm]t.south west) {Սա գրքի ամբողջ վերջաբանն է. ուրիշ ոչինչ};
    \node at ([xshift=-3.3cm, yshift=-0.2cm]current page.east) {\mascot[4.4cm]{smug}};
  \end{tikzpicture}
\end{frame}

% One-frame look back over the whole series
\begin{frame}{Ամբողջ շարքը՝ մեկ էջում}
  \centering
  \begin{tikzpicture}[ch/.style={draw=plum, line width=0.8pt, rounded corners=2mm, text width=4.25cm, minimum height=1.32cm,
      align=left, inner sep=3pt, font=\scriptsize, text=plum}]
    \node[ch, fill=butter!45!cream] at (0,1.5) {\emoji{star}\ {\blackfont\footnotesize Ներածություն}\\ Ուղեղն ապշեցուցիչ է, բայց կատարյալ չէ};
    \node[ch, fill=mint!40!cardwhite] at (4.6,1.5) {\emoji{rollercoaster}\ {\blackfont\footnotesize 1. Մարմնի կառավարում}\\ Ծովախտ, աղանդեր, քուն. հին ուղեղն ընդդեմ նորի};
    \node[ch, fill=lavender!45!cardwhite] at (9.2,1.5) {\emoji{thought}\ {\blackfont\footnotesize 2. Հիշողություն}\\ Եսասեր է ու ամեն անգամ խմբագրվում է};
    \node[ch, fill=peach!45!cardwhite] at (0,0) {\emoji{scream}\ {\blackfont\footnotesize 3. Վախ}\\ Վախենում ենք նաև անհավանական ու անվնաս բաներից};
    \node[ch, fill=butter!45!cream] at (4.6,0) {\emoji{graduation}\ {\blackfont\footnotesize 4. Խելք}\\ Դժվար է չափել, ու խելացիներն էլ են հիմարություններ անում};
    \node[ch, fill=mint!40!cardwhite] at (9.2,0) {\emoji{eye}\ {\blackfont\footnotesize 5. Ընկալում}\\ Ընկալածը ուղեղի կռահումն է. դեմքեր տոստի վրա};
    \node[ch, fill=lavender!45!cardwhite] at (0,-1.5) {\emoji{masks}\ {\blackfont\footnotesize 6. Անհատականություն}\\ Թեստերը թույլ են, զայրույթն ու հումորը՝ խճճված};
    \node[ch, fill=peach!45!cardwhite] at (4.6,-1.5) {\emoji{hug}\ {\blackfont\footnotesize 7. Ուրիշ մարդիկ}\\ Խմբեր, սեր, կոտրված սիրտ ու անարդարություն};
    \node[ch, fill=brainpink!40!cardwhite, line width=1.4pt] at (9.2,-1.5) {\emoji{brain}\ {\blackfont\footnotesize 8. Երբ ուղեղը «փչանում» է}\\ Հոգեկան խնդիրներն իրական են ու կարեկցանքի արժանի};
  \end{tikzpicture}\par\smallskip
  \onslide<2->{{\small \emoji{speech}\ Չաթում. ո՞ր գլուխն էր քո \textbf{սիրելին}։ Գրիր համարը}\par}
\end{frame}

% Goodbye
\begin{frame}{Շնորհակալություն}
  \begin{columns}[c]
    \begin{column}{0.58\textwidth}
      {\blackfont\large 7 վեբինար, 8 գլուխ, 1 ապուշ ուղեղ}\par\medskip
      \begin{card}[colback=mint!30!cardwhite]
        Հիմա գիտես, թե ինչպես է ուղեղդ քեզ խաբում, վախեցնում ու երբեմն «փչանում»։ Գուցե մի քիչ ավելի \textbf{բարի} լինես նրա ու ուրիշների ուղեղների նկատմամբ \emoji{heart}
      \end{card}
      \vspace{2mm}
      {\small Շնորհակալություն, որ մեզ հետ էիր ամբողջ շարքի ընթացքում \emoji{wavehand}}
    \end{column}
    \begin{column}{0.36\textwidth}
      \centering\mascot[3.4cm]{happy}
    \end{column}
  \end{columns}
\end{frame}

\begin{frame}[plain]
  \begin{tikzpicture}[remember picture, overlay]
    \fill[lavender!60!cream] (current page.south west) rectangle (current page.north east);
    \fill[pattern={Dots[distance=4.4mm, radius=0.28mm]}, pattern color=plum!10!lavender!60!cream]
      (current page.south west) rectangle (current page.north east);
    \node at ([xshift=-3.2cm]current page.center) {\mascot[4.4cm]{happy}};
    \node[anchor=west, text=plum] at ([xshift=0.2cm, yshift=0.6cm]current page.center)
      {\titlefont\fontsize{40}{42}\selectfont Հարցեր};
    \node[anchor=west, text=plum, font=\mediumfont\large, align=left] at ([xshift=0.2cm, yshift=-1.0cm]current page.center)
      {Գրիր չաթում կամ\\ միացրու խոսափողը \emoji{megaphone}};
    \node[anchor=west, text=plum!70!cream, font=\small] at ([xshift=0.2cm, yshift=-2.1cm]current page.center)
      {Դին Բըրնեթ, «The Idiot Brain», 2016};
  \end{tikzpicture}
\end{frame}
